%source mathadata
\textbf{Correction de l'exercice 3 :}
		\begin{enumerate}
			\item \textbf{Équation de la droite $(AB)$ :}

			      Pour trouver l'équation réduite, on cherche l'équation de la forme $y = mx + p$.

			      \begin{itemize}
				      \item Le coefficient directeur $m$ (ou pente) se calcule avec la formule :
				            $$m = \frac{y_B - y_A}{x_B - x_A} = \frac{0 - 3}{4 - 0} = \frac{-3}{4}$$

				      \item Pour trouver l'ordonnée à l'origine $p$, on utilise le point $A(0;3)$. Comme $A$ appartient à la droite, ses coordonnées vérifient l'équation :
				            $$y_A = m \times x_A + p$$
				            $$3 = -\frac{3}{4} \times 0 + p$$
				            $$p = 3$$
			      \end{itemize}

			      Donc l'équation de la droite $(AB)$ est : $\boxed{y = -\frac{3}{4}x + 3}$

			\item \textbf{Position du point $C(2;4)$ par rapport à la droite :}
			      \begin{enumerate}
				      \item \textbf{Coordonnées du point $E$ :}

				            Le point $E$ a la même abscisse que $C$, donc $x_E = 2$. Comme $E$ appartient à la droite $(AB)$, son ordonnée vérifie l'équation :
				            $$y_E = -\frac{3}{4} \times 2 + 3 = -\frac{6}{4} + 3 = -\frac{3}{2} + 3 = \frac{3}{2} = 1,5$$

				            Donc $\boxed{E(2;1,5)}$

				      \item \textbf{Comparaison des ordonnées :}

				            L'ordonnée de $C$ est $y_C = 4$.

				            L'ordonnée de $E$ est $y_E = 1,5$.

				            On constate que $y_C = 4 > 1,5 = y_E$

				      \item \textbf{Conclusion :}

				            Comme l'ordonnée de $C$ est plus grande que celle de $E$ (qui est sur la droite à la même abscisse), le point $C$ est \boxed{\text{au-dessus}} de la droite $(AB)$.
			      \end{enumerate}

			\item \textbf{Méthode générale (Bonus) :}

			      Pour déterminer si un point $M(x_M;y_M)$ est au-dessus ou en dessous d'une droite d'équation $y = mx + p$ :
			      \begin{enumerate}
				      \item On calcule $y_{\text{droite}} = m \times x_M + p$ (ordonnée du point de la droite ayant la même abscisse que $M$).
				      \item On compare $y_M$ avec $y_{\text{droite}}$ :
				            \begin{itemize}
					            \item Si $y_M > y_{\text{droite}}$, alors $M$ est au-dessus de la droite.
					            \item Si $y_M < y_{\text{droite}}$, alors $M$ est en dessous de la droite.
					            \item Si $y_M = y_{\text{droite}}$, alors $M$ appartient à la droite.
				            \end{itemize}
			      \end{enumerate}
		\end{enumerate}
